\section{More about lengths and simples}

Let us continue studying lengths of permutations.

% ============================================================
% Proposition prop.perm.len.inv
% ============================================================

\begin{proposition}
\label{prop.perm.len.inv}
\lean{Equiv.Perm.length_inv}
\leantarget
\leanok
Let $n \in \mathbb{N}$ and $\sigma \in S_n$. Then,
$\ell(\sigma^{-1}) = \ell(\sigma)$.
\end{proposition}

\begin{proof}
\leanok
Recall that $\ell(\sigma)$ is the number of inversions of $\sigma$,
while $\ell(\sigma^{-1})$ is the number of inversions of $\sigma^{-1}$.

An inversion of $\sigma$ is a pair $(i,j) \in [n] \times [n]$ such that
$i < j$ and $\sigma(i) > \sigma(j)$. Likewise, an inversion of $\sigma^{-1}$
is a pair $(u,v) \in [n] \times [n]$ such that $u < v$ and
$\sigma^{-1}(u) > \sigma^{-1}(v)$.

Thus, if $(i,j)$ is an inversion of $\sigma$, then
$(\sigma(j), \sigma(i))$ is an inversion of $\sigma^{-1}$.
Hence, we obtain a map
\begin{align*}
\{\text{inversions of } \sigma\} &\to \{\text{inversions of } \sigma^{-1}\}, \\
(i,j) &\mapsto (\sigma(j), \sigma(i)).
\end{align*}
This map is bijective (indeed, it has an inverse map, which sends
each $(u,v) \in \{\text{inversions of } \sigma^{-1}\}$ to
$(\sigma^{-1}(v), \sigma^{-1}(u))$). Thus, the bijection principle yields
\[
(\text{number of inversions of } \sigma) = (\text{number of inversions of } \sigma^{-1}).
\]
In other words, $\ell(\sigma) = \ell(\sigma^{-1})$.
(See \cite[Exercise 5.2 \textbf{(f)}]{detnotes} for details.)
\end{proof}

% ============================================================
% Helper: inversionsBijection (Lean-only)
% ============================================================

\begin{definition}
\label{def.perm.inversionsBijection}
\lean{Equiv.Perm.inversionsBijection}
\leanhelper
\leanok
The explicit bijection between inversions of $\sigma$ and inversions of $\sigma^{-1}$,
sending $(i,j)$ to $(\sigma(j), \sigma(i))$.
\end{definition}

% ============================================================
% Lemma lem.perm.len.ssl
% ============================================================

\begin{lemma}[Single swap lemma]
\label{lem.perm.len.ssl}
\lean{Equiv.Perm.length_mul_swap_right, Equiv.Perm.length_mul_swap_left}
\leantarget
\leanok
Let $n \in \mathbb{N}$, $\sigma \in S_n$ and $k \in [n-1]$. Then:
\medskip

\textbf{(a)} We have
\[
\ell(\sigma s_k) =
\begin{cases}
\ell(\sigma) + 1, & \text{if } \sigma(k) < \sigma(k+1); \\
\ell(\sigma) - 1, & \text{if } \sigma(k) > \sigma(k+1).
\end{cases}
\]

\textbf{(b)} We have
\[
\ell(s_k \sigma) =
\begin{cases}
\ell(\sigma) + 1, & \text{if } \sigma^{-1}(k) < \sigma^{-1}(k+1); \\
\ell(\sigma) - 1, & \text{if } \sigma^{-1}(k) > \sigma^{-1}(k+1).
\end{cases}
\]
\end{lemma}

\begin{proof}
This combines parts \textbf{(a)} and \textbf{(b)}.
\end{proof}

\begin{theorem}
\label{lem.perm.len.ssl.a}
\lean{Equiv.Perm.length_mul_swap_right}
\leanhelper
\leanok
Part \textbf{(a)} of Lemma~\ref{lem.perm.len.ssl}: When multiplying a permutation $\sigma$
on the right by the simple transposition $s_k$, we have
$\ell(\sigma s_k) = \ell(\sigma) + 1$ if $\sigma(k) < \sigma(k+1)$,
and $\ell(\sigma s_k) = \ell(\sigma) - 1$ if $\sigma(k) > \sigma(k+1)$.
\end{theorem}

\begin{proof}
\leanok
The proof partitions the inversions of $\sigma$ and $\sigma s_k$ into two kinds.
Call an inversion $(i,j)$ of a permutation $\tau$ the \emph{key pair}
if $(i,j) = (k,k+1)$; all other inversions are \emph{non-key}.

\emph{Key observation:} There is a bijection between non-key inversions
of $\sigma$ and non-key inversions of $\sigma s_k$, given by applying
the swap $k \leftrightarrow k+1$ to both coordinates.
This bijection preserves the strict ordering of the pair coordinates
because there is no element strictly between $k$ and $k+1$.

Thus
\[
(\text{non-key inversions of } \sigma s_k) = (\text{non-key inversions of } \sigma).
\]

For the key pair $(k,k+1)$: it is an inversion of $\sigma$ if and only if
$\sigma(k+1) < \sigma(k)$, and it is an inversion of $\sigma s_k$ if and only if
$\sigma(k) < \sigma(k+1)$.
So the key pair's inversion status \emph{flips} when we pass from $\sigma$ to $\sigma s_k$.

Combining these two observations: the number of inversions changes by exactly $\pm 1$,
depending on whether the key pair is added or removed.
\end{proof}

\begin{theorem}
\label{lem.perm.len.ssl.b}
\lean{Equiv.Perm.length_mul_swap_left}
\leanhelper
\leanok
Part \textbf{(b)} of Lemma~\ref{lem.perm.len.ssl}: When multiplying a permutation $\sigma$
on the left by the simple transposition $s_k$, we have
$\ell(s_k \sigma) = \ell(\sigma) + 1$ if $\sigma^{-1}(k) < \sigma^{-1}(k+1)$,
and $\ell(s_k \sigma) = \ell(\sigma) - 1$ if $\sigma^{-1}(k) > \sigma^{-1}(k+1)$.
\end{theorem}

\begin{proof}
\leanok
We use the identity $s_k \sigma = (\sigma^{-1} s_k)^{-1}$
(since $s_k$ is its own inverse)
and then apply Proposition~\ref{prop.perm.len.inv} and part~\textbf{(a)}.
\end{proof}

% ============================================================
% Proposition prop.perm.lisitij
% ============================================================

\begin{proposition}
\label{prop.perm.lisitij}
\lean{Equiv.Perm.length_mul_transposition_compare, Equiv.Perm.length_mul_transposition}
\leantarget
\leanok
Let $n \in \mathbb{N}$ and $\sigma \in S_n$. Let $i$ and $j$ be
two elements of $[n]$ such that $i < j$. Then:
\medskip

\textbf{(a)} We have $\ell(\sigma t_{i,j}) < \ell(\sigma)$ if $\sigma(i) > \sigma(j)$.
We have $\ell(\sigma t_{i,j}) > \ell(\sigma)$ if $\sigma(i) < \sigma(j)$.
\medskip

\textbf{(b)} We have
\[
\ell(\sigma t_{i,j}) =
\begin{cases}
\ell(\sigma) - 2|Q| - 1, & \text{if } \sigma(i) > \sigma(j); \\
\ell(\sigma) + 2|R| + 1, & \text{if } \sigma(i) < \sigma(j),
\end{cases}
\]
where
\begin{align*}
Q &= \{k \in \{i+1, i+2, \ldots, j-1\} \mid \sigma(i) > \sigma(k) > \sigma(j)\}
\quad \text{and} \\
R &= \{k \in \{i+1, i+2, \ldots, j-1\} \mid \sigma(i) < \sigma(k) < \sigma(j)\}.
\end{align*}
\end{proposition}

\begin{proof}
This combines parts \textbf{(a)} and \textbf{(b)}.
\end{proof}

\begin{definition}
\label{def.perm.setQ}
\lean{Equiv.Perm.setQ}
\leanhelper
\leanok
The set $Q$ from Proposition~\ref{prop.perm.lisitij}:
$Q = \{k \in \{i+1, \ldots, j-1\} \mid \sigma(j) < \sigma(k) < \sigma(i)\}$.
\end{definition}

\begin{definition}
\label{def.perm.setR}
\lean{Equiv.Perm.setR}
\leanhelper
\leanok
The set $R$ from Proposition~\ref{prop.perm.lisitij}:
$R = \{k \in \{i+1, \ldots, j-1\} \mid \sigma(i) < \sigma(k) < \sigma(j)\}$.
\end{definition}

\subsection{Helpers for Proposition~\ref{prop.perm.lisitij}}

\begin{definition}
\label{def.perm.setA}
\lean{Equiv.Perm.setA}
\leanhelper
\leanok
The set $A$ from the transposition length proof:
$A = \{a < i \mid \sigma(j) < \sigma(a) < \sigma(i)\}$.
These elements produce paired lost and gained inversions:
$(a, j)$ is lost and $(a, i)$ is gained.
\end{definition}

\begin{definition}
\label{def.perm.setB}
\lean{Equiv.Perm.setB}
\leanhelper
\leanok
The set $B$ from the transposition length proof:
$B = \{b > j \mid \sigma(j) < \sigma(b) < \sigma(i)\}$.
These elements produce paired lost and gained inversions:
$(i, b)$ is lost and $(j, b)$ is gained.
\end{definition}

\begin{lemma}
\label{lem.perm.setQ_swap_eq_setR}
\lean{Equiv.Perm.setQ_swap_eq_setR}
\leanhelper
\leanok
The set $Q$ for $\sigma \cdot t_{i,j}$ equals the set $R$ for $\sigma$.
This key symmetry allows reducing the case $\sigma(i) < \sigma(j)$
to the case $\sigma(i) > \sigma(j)$ in the transposition length formula.
\end{lemma}

\begin{proof}
\leanok
For $k$ between $i$ and $j$ with $k \neq i, j$, the swap $t_{i,j}$ fixes $k$,
so $(\sigma \cdot t_{i,j})(k) = \sigma(k)$. The conditions
$(\sigma t_{i,j})(j) < \sigma(k) < (\sigma t_{i,j})(i)$ become
$\sigma(i) < \sigma(k) < \sigma(j)$, which defines $R$.
\end{proof}

\begin{lemma}
\label{lem.perm.ij_lost}
\lean{Equiv.Perm.ij_lost}
\leanhelper
\leanok
When $\sigma(j) < \sigma(i)$, the pair $(i,j)$ itself is a lost inversion:
$(i,j) \in \operatorname{inv}(\sigma) \setminus \operatorname{inv}(\sigma \cdot t_{i,j})$.
\end{lemma}

\begin{proof}
\leanok
Since $(\sigma t_{i,j})(i) = \sigma(j) < \sigma(i) = (\sigma t_{i,j})(j)$,
the pair $(i,j)$ is no longer an inversion after applying $t_{i,j}$.
\end{proof}

\begin{lemma}
\label{lem.perm.ik_lost_for_k_in_Q}
\lean{Equiv.Perm.ik_lost_for_k_in_Q}
\leanhelper
\leanok
For $k \in Q$: the pair $(i, k)$ is a lost inversion.
Before the swap, $\sigma(i) > \sigma(k)$ so $(i,k)$ is an inversion.
After the swap, $(\sigma t_{i,j})(i) = \sigma(j) < \sigma(k)$, so it is not.
\end{lemma}

\begin{proof}
\leanok
Direct verification using the definition of $Q$ and the swap.
\end{proof}

\begin{lemma}
\label{lem.perm.kj_lost_for_k_in_Q}
\lean{Equiv.Perm.kj_lost_for_k_in_Q}
\leanhelper
\leanok
For $k \in Q$: the pair $(k, j)$ is a lost inversion.
Before the swap, $\sigma(k) > \sigma(j)$ so $(k,j)$ is an inversion.
After the swap, $(\sigma t_{i,j})(j) = \sigma(i) > \sigma(k)$, so it is not.
\end{lemma}

\begin{proof}
\leanok
Direct verification using the definition of $Q$ and the swap.
\end{proof}

\begin{lemma}
\label{lem.perm.aj_lost_for_a_in_A}
\lean{Equiv.Perm.aj_lost_for_a_in_A}
\leanhelper
\leanok
For $a \in A$: the pair $(a, j)$ is a lost inversion.
Before the swap, $\sigma(a) > \sigma(j)$. After the swap,
$(\sigma t_{i,j})(j) = \sigma(i) > \sigma(a)$, so $(a,j)$ is no longer an inversion.
\end{lemma}

\begin{proof}
\leanok
Direct verification.
\end{proof}

\begin{lemma}
\label{lem.perm.ib_lost_for_b_in_B}
\lean{Equiv.Perm.ib_lost_for_b_in_B}
\leanhelper
\leanok
For $b \in B$: the pair $(i, b)$ is a lost inversion.
Before the swap, $\sigma(i) > \sigma(b)$. After the swap,
$(\sigma t_{i,j})(i) = \sigma(j) < \sigma(b)$, so $(i,b)$ is no longer an inversion.
\end{lemma}

\begin{proof}
\leanok
Direct verification.
\end{proof}

\begin{lemma}
\label{lem.perm.ai_gained_for_a_in_A}
\lean{Equiv.Perm.ai_gained_for_a_in_A}
\leanhelper
\leanok
For $a \in A$: the pair $(a, i)$ is a gained inversion.
Before the swap, $\sigma(a) < \sigma(i)$ so $(a,i)$ is not an inversion.
After the swap, $(\sigma t_{i,j})(i) = \sigma(j) < \sigma(a)$, so it becomes one.
\end{lemma}

\begin{proof}
\leanok
Direct verification.
\end{proof}

\begin{lemma}
\label{lem.perm.jb_gained_for_b_in_B}
\lean{Equiv.Perm.jb_gained_for_b_in_B}
\leanhelper
\leanok
For $b \in B$: the pair $(j, b)$ is a gained inversion.
Before the swap, $\sigma(j) < \sigma(b)$ so $(j,b)$ is not an inversion.
After the swap, $(\sigma t_{i,j})(j) = \sigma(i) > \sigma(b)$, so it becomes one.
\end{lemma}

\begin{proof}
\leanok
Direct verification.
\end{proof}

\begin{lemma}
\label{lem.perm.lost_inversion_complete}
\lean{Equiv.Perm.lost_inversion_complete}
\leanhelper
\leanok
Completeness of the lost inversion classification:
every lost inversion $(a,b) \in \operatorname{inv}(\sigma) \setminus \operatorname{inv}(\sigma \cdot t_{i,j})$
(with $\sigma(j) < \sigma(i)$) is one of five types:
\begin{enumerate}
\item $(a,b) = (i,j)$,
\item $a \in A$ and $b = j$,
\item $a = i$ and $b \in B$,
\item $a = i$ and $b \in Q$,
\item $a \in Q$ and $b = j$.
\end{enumerate}
\end{lemma}

\begin{proof}
\leanok
By case analysis on whether $a$ or $b$ equals $i$ or $j$,
and using the fact that the swap only changes the values at positions $i$ and $j$.
\end{proof}

\begin{lemma}
\label{lem.perm.gained_inversion_complete}
\lean{Equiv.Perm.gained_inversion_complete}
\leanhelper
\leanok
Completeness of the gained inversion classification:
every gained inversion $(a,b) \in \operatorname{inv}(\sigma \cdot t_{i,j}) \setminus \operatorname{inv}(\sigma)$
(with $\sigma(j) < \sigma(i)$) is one of two types:
\begin{enumerate}
\item $a \in A$ and $b = i$,
\item $a = j$ and $b \in B$.
\end{enumerate}
\end{lemma}

\begin{proof}
By case analysis, similar to the lost inversion completeness.
\end{proof}

\begin{lemma}
\label{lem.perm.lost-minus-gained}
\lean{Equiv.Perm.lost_minus_gained_eq}
\leanhelper
\leanok
Let $\sigma \in S_n$ and $i < j$ with $\sigma(j) < \sigma(i)$.
Then the number of ``lost'' inversions (inversions of $\sigma$ that are not
inversions of $\sigma \cdot t_{i,j}$) exceeds the number of ``gained'' inversions
(inversions of $\sigma \cdot t_{i,j}$ that are not inversions of $\sigma$)
by exactly $2|Q| + 1$.

The proof partitions lost inversions into five disjoint types:
the pair $(i,j)$ itself; pairs $(a,j)$ and $(i,b)$ for $a \in A$, $b \in B$;
and pairs $(i,k)$ and $(k,j)$ for $k \in Q$.
Gained inversions are pairs $(a,i)$ for $a \in A$ and $(j,b)$ for $b \in B$.
The sets $A$ and $B$ cancel, leaving $2|Q|+1$.
\end{lemma}

\begin{proof}
\leanok
By explicit enumeration and counting of the five types of lost inversions
and two types of gained inversions.
\end{proof}

\begin{theorem}
\label{prop.perm.lisitij.b}
\lean{Equiv.Perm.length_mul_transposition}
\leanhelper
\leanok
Part \textbf{(b)} of Proposition~\ref{prop.perm.lisitij}: The exact formula for length change
when multiplying by an arbitrary transposition $t_{i,j}$.
\end{theorem}

\begin{proof}
\leanok
\textbf{Case $\sigma(i) > \sigma(j)$:}
This is \cite[Exercise 5.20]{detnotes}.
Apply Lemma~\ref{lem.perm.lost-minus-gained} directly.
The symmetric difference formula
$|\mathrm{inv}(\sigma t_{i,j})| + |\mathrm{lost}| = |\mathrm{inv}(\sigma)| + |\mathrm{gained}|$
together with $|\mathrm{lost}| = |\mathrm{gained}| + 2|Q| + 1$
gives $\ell(\sigma t_{i,j}) = \ell(\sigma) - 2|Q| - 1$.

\textbf{Case $\sigma(i) < \sigma(j)$:}
Apply the previous case to $\sigma t_{i,j}$ instead of $\sigma$.
Since $(\sigma t_{i,j})(i) = \sigma(j) > \sigma(i) = (\sigma t_{i,j})(j)$,
and $\sigma t_{i,j} t_{i,j} = \sigma$,
and the sets $Q$ and $R$ trade places when we replace $\sigma$ by $\sigma t_{i,j}$,
we obtain $\ell(\sigma) = \ell(\sigma t_{i,j}) - 2|R| - 1$, hence
$\ell(\sigma t_{i,j}) = \ell(\sigma) + 2|R| + 1$.
\end{proof}

\begin{theorem}
\label{prop.perm.lisitij.a}
\lean{Equiv.Perm.length_mul_transposition_compare}
\leanhelper
\leanok
Part \textbf{(a)} of Proposition~\ref{prop.perm.lisitij}: The length decreases
if $(i,j)$ was an inversion, and increases if it was not.
\end{theorem}

\begin{proof}
\leanok
This follows immediately from part \textbf{(b)}: when $\sigma(i) > \sigma(j)$,
the formula gives $\ell(\sigma t_{i,j}) = \ell(\sigma) - 2|Q| - 1 < \ell(\sigma)$,
and when $\sigma(i) < \sigma(j)$, the formula gives
$\ell(\sigma t_{i,j}) = \ell(\sigma) + 2|R| + 1 > \ell(\sigma)$.
\end{proof}

% ============================================================
% Convention on simple transpositions
% ============================================================

\begin{convention}
\label{conv.perm.simple}
\lean{AlgebraicCombinatorics.simpleTransposition}
\leanok
A \emph{simple transposition} in $S_n$ means one of the $n-1$ transpositions
$s_0, s_1, \ldots, s_{n-2}$. We shall occasionally abbreviate
``simple transposition'' as ``simple''.
\end{convention}

\begin{definition}
\label{def.perm.Word}
\lean{Equiv.Perm.Word}
\leanhelper
\leanok
A \emph{word} (of length $q$) in $S_n$ is a finite sequence $(k_1, k_2, \ldots, k_q)$
where each $k_i \in \{0, 1, \ldots, n-2\}$ is an index of a simple transposition.
\end{definition}

\begin{definition}
\label{def.perm.wordProd}
\lean{Equiv.Perm.wordProd}
\leanhelper
\leanok
The \emph{product} of a word $w = (k_1, \ldots, k_q)$ is the
permutation $s_{k_1} s_{k_2} \cdots s_{k_q}$.
\end{definition}

\begin{definition}
\label{def.perm.IsReduced}
\lean{Equiv.Perm.IsReduced}
\leanhelper
\leanok
A word $w = (k_1, \ldots, k_q)$ is \emph{reduced} if its length $q$ equals
$\ell(s_{k_1} \cdots s_{k_q})$.
\end{definition}

% ============================================================
% Theorem thm.perm.len.redword1
% ============================================================

\begin{theorem}[First reduced word theorem for the symmetric group]
\label{thm.perm.len.redword1}
\lean{Equiv.Perm.exists_reduced_word, Equiv.Perm.length_le_word_length}
\leantarget
\leanok
Let $n \in \mathbb{N}$ and $\sigma \in S_n$. Then:
\medskip

\textbf{(a)} We can write $\sigma$ as a composition (i.e., product) of
$\ell(\sigma)$ simples.
\medskip

\textbf{(b)} The number $\ell(\sigma)$ is the smallest $p \in \mathbb{N}$
such that we can write $\sigma$ as a composition of $p$ simples.
\medskip

\textbf{Keep in mind:} The composition of $0$ simples is $\mathrm{id}$,
since $\mathrm{id}$ is the neutral element of the group $S_n$.
\end{theorem}

\begin{proof}
This combines parts \textbf{(a)} and \textbf{(b)}.
\end{proof}

\begin{theorem}
\label{thm.perm.len.redword1.a}
\lean{Equiv.Perm.exists_reduced_word}
\leanhelper
\leanok
Part \textbf{(a)} of Theorem~\ref{thm.perm.len.redword1}: Every permutation $\sigma \in S_n$
can be written as a composition of $\ell(\sigma)$ simple transpositions.
Equivalently, there exists a reduced word $w$ whose product $s_{k_1} \cdots s_{k_q} = \sigma$.
\end{theorem}

\begin{proof}
\leanok
(See \cite[Exercise 5.2 \textbf{(e)}]{detnotes} for details.)

We proceed by strong induction on $\ell(\sigma)$.

\emph{Base case:} If $\ell(\sigma) = 0$, then $\sigma = \mathrm{id}$
(a permutation with no inversions is the identity), so the empty word works.

\emph{Induction step:} If $\ell(\sigma) = h + 1 > 0$, then $\sigma$ has at least
one inversion, so there exists $k \in [n-1]$ with $\sigma(k) > \sigma(k+1)$.
By Lemma~\ref{lem.perm.len.ssl}~\textbf{(a)},
$\ell(\sigma s_k) = \ell(\sigma) - 1 = h$.
By the induction hypothesis, $\sigma s_k$ has a reduced word $w$ of length $h$.
Then $w \mathbin{+\!+} [k]$ is a reduced word of length $h+1$ for
$\sigma s_k \cdot s_k = \sigma$.
\end{proof}

\begin{theorem}
\label{thm.perm.len.redword1.b}
\lean{Equiv.Perm.length_le_word_length}
\leanhelper
\leanok
Part \textbf{(b)} of Theorem~\ref{thm.perm.len.redword1}:
For any word $w = (k_1, \ldots, k_q)$, we have $\ell(s_{k_1} \cdots s_{k_q}) \leq q$.
In other words, $\ell(\sigma)$ is the minimum word length.
\end{theorem}

\begin{proof}
\leanok
(See \cite[Exercise 5.2 \textbf{(g)}]{detnotes} for details.)

We prove $\ell(s_{k_1} s_{k_2} \cdots s_{k_g}) \leq g$ by induction on $g$.
The base case $g = 0$ gives $\ell(\mathrm{id}) = 0$.
For the induction step, Lemma~\ref{lem.perm.len.ssl}~\textbf{(b)} gives
$\ell(s_k \cdot \tau) \leq \ell(\tau) + 1$ for any $\tau$ and $k$,
so $\ell(s_{k_1} \cdots s_{k_g}) \leq \ell(s_{k_2} \cdots s_{k_g}) + 1 \leq (g-1) + 1 = g$.
\end{proof}

\begin{theorem}
\label{thm.perm.len.eq-min}
\lean{Equiv.Perm.length_eq_min_word_length}
\leanhelper
\leanok
Combining parts \textbf{(a)} and \textbf{(b)}: $\ell(\sigma)$ is exactly the
minimum word length for $\sigma$.
\end{theorem}

\begin{proof}
\leanok
Immediate from parts \textbf{(a)} and \textbf{(b)}.
\end{proof}

% ============================================================
% Corollary cor.perm.red.sigtau
% ============================================================

\begin{corollary}
\label{cor.perm.red.sigtau}
\lean{Equiv.Perm.length_mul_mod_two, Equiv.Perm.length_mul_le, Equiv.Perm.word_length_ge_and_parity}
\leantarget
\leanok
Let $n \in \mathbb{N}$.
\medskip

\textbf{(a)} We have $\ell(\sigma\tau) \equiv \ell(\sigma) + \ell(\tau) \pmod{2}$
for all $\sigma, \tau \in S_n$.
\medskip

\textbf{(b)} We have $\ell(\sigma\tau) \leq \ell(\sigma) + \ell(\tau)$
for all $\sigma, \tau \in S_n$.
\medskip

\textbf{(c)} Let $k_1, k_2, \ldots, k_q \in [n-1]$, and let
$\sigma = s_{k_1} s_{k_2} \cdots s_{k_q}$. Then $q \geq \ell(\sigma)$
and $q \equiv \ell(\sigma) \pmod{2}$.
\end{corollary}

\begin{proof}
This combines parts \textbf{(a)}, \textbf{(b)}, and \textbf{(c)}.
\end{proof}

\begin{theorem}
\label{cor.perm.red.sigtau.a}
\lean{Equiv.Perm.length_mul_mod_two}
\leanhelper
\leanok
Part \textbf{(a)} of Corollary~\ref{cor.perm.red.sigtau}: The length of a product
has the same parity as the sum of lengths.
\end{theorem}

\begin{proof}
\leanok
(See \cite[Exercise 5.2 \textbf{(b)}]{detnotes} for details.)

Write $\tau = s_{k_1} s_{k_2} \cdots s_{k_q}$ with $q = \ell(\tau)$
(using Theorem~\ref{thm.perm.len.redword1}~\textbf{(a)}). Then
\begin{align*}
\ell(\sigma\tau) &= \ell(\sigma s_{k_1} s_{k_2} \cdots s_{k_q}) \\
&\equiv \ell(\sigma s_{k_1} \cdots s_{k_{q-1}}) + 1 \\
&\equiv \cdots \\
&\equiv \ell(\sigma) + q = \ell(\sigma) + \ell(\tau) \pmod{2}.
\end{align*}
\end{proof}

\begin{theorem}
\label{cor.perm.red.sigtau.b}
\lean{Equiv.Perm.length_mul_le}
\leanhelper
\leanok
Part \textbf{(b)} of Corollary~\ref{cor.perm.red.sigtau}: The triangle inequality
$\ell(\sigma\tau) \leq \ell(\sigma) + \ell(\tau)$.
\end{theorem}

\begin{proof}
\leanok
(See \cite[Exercise 5.2 \textbf{(c)}]{detnotes} for details.)

Get reduced words $w_1, w_2$ for $\sigma, \tau$. Then $w_1 \mathbin{+\!+} w_2$
represents $\sigma\tau$, so
$\ell(\sigma\tau) \leq |w_1 \mathbin{+\!+} w_2| = |w_1| + |w_2| = \ell(\sigma) + \ell(\tau)$
by Theorem~\ref{thm.perm.len.redword1}~\textbf{(b)}.
\end{proof}

\begin{theorem}
\label{cor.perm.red.sigtau.c}
\lean{Equiv.Perm.word_length_ge_and_parity}
\leanhelper
\leanok
Part \textbf{(c)} of Corollary~\ref{cor.perm.red.sigtau}: For any word representing $\sigma$,
the word length is at least $\ell(\sigma)$ and has the same parity as $\ell(\sigma)$.
\end{theorem}

\begin{proof}
\leanok
The inequality $q \geq \ell(\sigma)$ is
Theorem~\ref{thm.perm.len.redword1}~\textbf{(b)}.
The parity claim $q \equiv \ell(\sigma) \pmod{2}$ is proved by
the same telescoping argument as part~\textbf{(a)}, starting from
$\ell(\sigma) = \ell(s_{k_1} \cdots s_{k_q})$ and peeling off one
simple at a time.
\end{proof}

% ============================================================
% Corollary cor.perm.generated
% ============================================================

\begin{corollary}
\label{cor.perm.generated}
\lean{Equiv.Perm.generated_by_simples}
\leantarget
\leanok
Let $n \in \mathbb{N}$. Then, the group $S_n$ is generated by the simples
$s_0, s_1, \ldots, s_{n-2}$.
\end{corollary}

\begin{proof}
\leanok
This follows directly from Theorem~\ref{thm.perm.len.redword1}~\textbf{(a)}:
every $\sigma \in S_n$ can be written as a product of simple transpositions,
so every $\sigma$ lies in the subgroup generated by the simples.
\end{proof}

% ============================================================
% Proposition prop.perm.redword-lehmer
% ============================================================

\subsection{Rothe diagram and Lehmer entry properties}

\begin{definition}
\label{def.perm.inRotheDiagram}
\lean{Equiv.Perm.inRotheDiagram}
\leanhelper
\leanok
A cell $(i, j)$ is in the \emph{Rothe diagram} of $\sigma$ if
$j < \sigma(i)$ and $i < \sigma^{-1}(j)$.
These are the cells that are not ``hit'' by either the row or column
of a permutation matrix entry.
\end{definition}

\begin{lemma}
\label{lem.perm.rotheDiagram_card_eq_length}
\lean{Equiv.Perm.rotheDiagram_card_eq_length}
\leanhelper
\leanok
The number of cells in the Rothe diagram of $\sigma$ equals $\ell(\sigma)$.
The bijection sends $(i, j)$ in the Rothe diagram to the inversion $(i, \sigma^{-1}(j))$.
\end{lemma}

\begin{proof}
\leanok
The map $(i,j) \mapsto (i, \sigma^{-1}(j))$ is a bijection from Rothe diagram cells
to inversions: the condition $j < \sigma(i)$ becomes $\sigma(\sigma^{-1}(j)) < \sigma(i)$,
which is $\sigma(k) < \sigma(i)$ with $k = \sigma^{-1}(j) > i$.
\end{proof}

\begin{lemma}
\label{lem.perm.lehmerEntry_lt2}
\lean{Equiv.Perm.lehmerEntry_lt}
\leanhelper
\leanok
For $\sigma \in S_n$ and $i \in [n]$, we have $\ell_i(\sigma) < n - i$.
This strict bound holds because $\ell_i(\sigma)$ counts elements among
$\{i+1, \ldots, n-1\}$ where $\sigma(j) < \sigma(i)$,
and there are only $n - 1 - i$ such elements.
\end{lemma}

\begin{proof}
\leanok
The set counted by $\ell_i(\sigma)$ is a subset of $\{j : j > i\}$,
which has cardinality $n - 1 - i < n - i$.
\end{proof}

\begin{lemma}
\label{lem.perm.sigma_zero_eq_lehmerEntry}
\lean{Equiv.Perm.sigma_zero_eq_lehmerEntry}
\leanhelper
\leanok
For $\sigma \in S_n$ with $n > 0$, we have $\sigma(0) = \ell_0(\sigma)$.
This is because $\ell_0(\sigma)$ counts how many elements among
$\sigma(1), \ldots, \sigma(n-1)$ are less than $\sigma(0)$,
and since $\sigma$ is a bijection onto $\{0, \ldots, n-1\}$,
this count equals $\sigma(0)$.
\end{lemma}

\begin{proof}
The elements $j > 0$ with $\sigma(j) < \sigma(0)$ are in bijection
(via $j \mapsto \sigma(j)$) with elements $v < \sigma(0)$.
Since $\sigma$ is a bijection, there are exactly $\sigma(0)$ such elements.
\end{proof}

\begin{lemma}
\label{lem.perm.lehmerEntry_diff_iff_inversion}
\lean{Equiv.Perm.lehmerEntry_diff_iff_inversion}
\leanhelper
\leanok
For $\sigma \in S_n$ and $i$ with $i + 1 < n$, let $i' = i + 1$. Then
\[
\ell_i(\sigma) \geq \ell_{i'}(\sigma) + 1 \iff \sigma(i) > \sigma(i').
\]
This characterizes when consecutive Lehmer entries differ by at least 1,
which is crucial for determining when a Lehmer block shifts a position.
\end{lemma}

\begin{proof}
The Lehmer entry $\ell_i(\sigma)$ counts $j > i$ with $\sigma(j) < \sigma(i)$.
We split this count by whether $j = i'$ or $j > i'$:
$\ell_i(\sigma) = [[\sigma(i') < \sigma(i)]] + |\{j > i' : \sigma(j) < \sigma(i)\}|$.
The set $\{j > i' : \sigma(j) < \sigma(i)\}$ contains
$\{j > i' : \sigma(j) < \sigma(i')\}$ (counted by $\ell_{i'}(\sigma)$)
when $\sigma(i') < \sigma(i)$, from which the equivalence follows.
\end{proof}

\begin{lemma}
\label{lem.perm.eq_of_inversions_eq}
\lean{Equiv.Perm.eq_of_inversions_eq}
\leanhelper
\leanok
Two permutations $\sigma, \tau \in S_n$ with the same inversions are equal:
if for all $i < j$ we have $\sigma(j) < \sigma(i) \iff \tau(j) < \tau(i)$,
then $\sigma = \tau$.
\end{lemma}

\begin{proof}
\leanok
For each $i$, the value $\sigma(i)$ equals $|\{j : \sigma(j) < \sigma(i)\}|$
(since $\sigma$ is a bijection onto $\{0, \ldots, n-1\}$).
The hypothesis implies $\{j : \sigma(j) < \sigma(i)\} = \{j : \tau(j) < \tau(i)\}$
for each $i$, so $\sigma(i) = \tau(i)$.
\end{proof}

\subsection{The Lehmer code representation}

\begin{definition}
\label{def.perm.lehmerBlock}
\lean{Equiv.Perm.lehmerBlock}
\leanhelper
\leanok
For each $i \in [n]$, we define the \emph{Lehmer block} $a_i$ as
\[
a_i := s_{i'-1} s_{i'-2} \cdots s_i
\]
where $i' = i + \ell_i(\sigma)$ and $\ell_i(\sigma)$ is the Lehmer entry at position $i$.
If $\ell_i(\sigma) = 0$, then $a_i = \mathrm{id}$ (the empty product).
The product $a_i$ equals the cycle $(i', i'-1, \ldots, i)$.
\end{definition}

\begin{lemma}
\label{lem.perm.lehmerBlock_length}
\lean{Equiv.Perm.lehmerBlock_length}
\leanhelper
\leanok
The length of the Lehmer block $a_i$ (as a word) equals the Lehmer entry $\ell_i(\sigma)$.
\end{lemma}

\begin{proof}
\leanok
The block consists of $\ell_i(\sigma)$ simple transpositions by construction.
\end{proof}

\begin{definition}
\label{def.perm.canonicalReducedWord}
\lean{Equiv.Perm.canonicalReducedWord}
\leanhelper
\leanok
The \emph{canonical reduced word} for $\sigma$ is the concatenation
of Lehmer blocks:
\[
w(\sigma) = a_0 \mathbin{+\!+} a_1 \mathbin{+\!+} \cdots \mathbin{+\!+} a_{n-1}.
\]
\end{definition}

\begin{lemma}
\label{lem.perm.canonicalReducedWord-length}
\lean{Equiv.Perm.canonicalReducedWord_length}
\leanhelper
\leanok
The length of the canonical reduced word equals $\ell(\sigma)$.
This follows from $\sum_{i} \ell_i(\sigma) = \ell(\sigma)$.
\end{lemma}

\begin{proof}
\leanok
Each block $a_i$ has length $\ell_i(\sigma)$, and the total length is
$\sum_i \ell_i(\sigma) = \ell(\sigma)$ (since each inversion contributes
to exactly one Lehmer entry).
\end{proof}

\begin{lemma}
\label{lem.perm.sum-lehmer-eq-length}
\lean{Equiv.Perm.sum_lehmerEntry_eq_length}
\leanhelper
\leanok
The sum of Lehmer entries equals the length:
$\sum_{i \in [n]} \ell_i(\sigma) = \ell(\sigma)$.
\end{lemma}

\begin{proof}
\leanok
By a bijection between the set of pairs $\{(i, j) : i \in [n],\ j \text{ such that } (i,j) \text{ is an inversion}\}$
and the set of all inversions.
\end{proof}

\subsection{Helpers for Proposition~\ref{prop.perm.redword-lehmer}}

\begin{definition}
\label{def.perm.smallerBefore}
\lean{Equiv.Perm.smallerBefore}
\leanhelper
\leanok
For $\sigma \in S_n$ and $i \in [n]$, define
$s(\sigma, i) = |\{j < i : \sigma(j) < \sigma(i)\}|$.
This counts positions before $i$ with smaller values.
\end{definition}

\begin{definition}
\label{def.perm.largerBefore}
\lean{Equiv.Perm.largerBefore}
\leanhelper
\leanok
For $\sigma \in S_n$ and $i \in [n]$, define
$g(\sigma, i) = |\{j < i : \sigma(j) > \sigma(i)\}|$.
This counts inversions with $i$ as the second element.
\end{definition}

\begin{lemma}
\label{lem.perm.smallerBefore_add_largerBefore}
\lean{Equiv.Perm.smallerBefore_add_largerBefore}
\leanhelper
\leanok
For $\sigma \in S_n$ and $i \in [n]$:
\[
s(\sigma, i) + g(\sigma, i) = i.
\]
This is because the elements $j < i$ are partitioned into those with $\sigma(j) < \sigma(i)$
and those with $\sigma(j) > \sigma(i)$ (equality is impossible since $\sigma$ is injective).
\end{lemma}

\begin{proof}
The set $\{j : j < i\}$ has cardinality $i$ and splits into
$\{j < i : \sigma(j) < \sigma(i)\}$ and $\{j < i : \sigma(j) > \sigma(i)\}$,
which are disjoint (since $\sigma$ is injective, $\sigma(j) \neq \sigma(i)$ for $j \neq i$).
\end{proof}

\begin{lemma}
\label{lem.perm.sigma-eq-formula}
\lean{Equiv.Perm.sigma_eq_smallerBefore_add_lehmerEntry}
\leanhelper
\leanok
For each $i \in [n]$, we have
$\sigma(i) = s(\sigma, i) + \ell_i(\sigma)$.
\end{lemma}

\begin{proof}
\leanok
The elements $j$ with $\sigma(j) < \sigma(i)$ partition into those with $j < i$
(counted by $s(\sigma, i)$) and those with $j > i$
(counted by $\ell_i$). Their total count is $\sigma(i)$.
\end{proof}

\begin{lemma}
\label{lem.perm.final_position_formula}
\lean{Equiv.Perm.final_position_formula}
\leanhelper
\leanok
For $\sigma \in S_n$ and $p \in [n]$:
\[
p + \ell_p(\sigma) - g(\sigma, p) = \sigma(p).
\]
This key formula shows that the final position after applying all Lehmer blocks
to position $p$ equals $\sigma(p)$.
\end{lemma}

\begin{proof}
From $\sigma(p) = s(\sigma, p) + \ell_p(\sigma)$
and $p = s(\sigma, p) + g(\sigma, p)$,
we get $p + \ell_p(\sigma) - g(\sigma, p) = \sigma(p)$.
\end{proof}

\begin{lemma}
\label{lem.perm.lehmerEntry_range_upper}
\lean{Equiv.Perm.lehmerEntry_range_upper}
\leanhelper
\leanok
For $\sigma \in S_n$ and $i \in [n]$:
\[
i + \ell_i(\sigma) = \sigma(i) + g(\sigma, i).
\]
This gives the upper bound of block $i$'s range in terms of the permutation value
and the count of larger elements before position $i$.
\end{lemma}

\begin{proof}
Follows from the formulas $\sigma(i) = s(\sigma, i) + \ell_i(\sigma)$
and $i = s(\sigma, i) + g(\sigma, i)$.
\end{proof}

\begin{proposition}
\label{prop.perm.redword-lehmer}
\lean{Equiv.Perm.canonicalReducedWord_prod}
\leantarget
\leanok
Let $n \in \mathbb{N}$ and $\sigma \in S_n$.
For each $i \in [n]$, set
\[
a_i := s_{i'-1} s_{i'-2} \cdots s_i,
\]
where $i' = i + \ell_i(\sigma)$. Then $\sigma = a_0 a_1 \cdots a_{n-1}$.

In other words, the product of the canonical reduced word equals $\sigma$.
\end{proposition}

\begin{proof}
We refer to \cite[Exercise 5.21 parts \textbf{(b)} and \textbf{(c)}]{detnotes}
for a detailed proof.

We prove by extensionality: for each position $p$, we track where it ends up
after applying all blocks.

By tracking position $p$ through all blocks, the final
position is
\[
p + \ell_p(\sigma) - g(\sigma, p).
\]

Since $p = s(\sigma, p) + g(\sigma, p)$
and $\sigma(p) = s(\sigma, p) + \ell_p(\sigma)$
(Lemma~\ref{lem.perm.sigma-eq-formula}), we get
\[
p + \ell_p(\sigma) - g(\sigma, p) = \sigma(p).
\]
\end{proof}

\begin{theorem}
\label{thm.perm.canonicalReducedWord-isReduced}
\lean{Equiv.Perm.canonicalReducedWord_isReduced}
\leanhelper
\leanok
The canonical reduced word is indeed reduced: it is a word of length $\ell(\sigma)$
whose product is $\sigma$.
\end{theorem}

\begin{proof}
\leanok
Combine Proposition~\ref{prop.perm.redword-lehmer}
(the product of the canonical reduced word equals $\sigma$)
with the length equality from Lemma~\ref{lem.perm.canonicalReducedWord-length}
($|w(\sigma)| = \ell(\sigma)$).
\end{proof}
